\section{Билет 64. Цикл Карно. Приведенная теплота. Приведенная теплота для обратимого цикла Карно. Термодинамическое определение энтропии. Математическая запись второго начала термодинамики.}

\begin{definition}
Величина $\delta Q / T$, где $\delta Q$ --- элементарное количество теплоты, полученное или отданное системой при абсолютной температуре $T$, называется \textbf{приведённой теплотой}.
\end{definition}

\begin{theorem}[Приведённая теплота для цикла Карно]
Для идеальной тепловой машины, работающей по обратимому циклу Карно между температурами нагревателя $T_1$ и холодильника $T_2$, КПД равен:
\begin{equation}
    \eta_{\text{К}} = 1 - \frac{T_2}{T_1} = 1 - \frac{|Q_2|}{Q_1}.
\end{equation}
Отсюда следует соотношение:
\begin{equation}
    \frac{Q_1}{T_1} = \frac{|Q_2|}{T_2} \quad \Rightarrow \quad \frac{Q_1}{T_1} + \frac{Q_2}{T_2} = 0, \label{eq:carnot_reduced_heat}
\end{equation}
где $Q_2 < 0$ --- теплота, отданная холодильнику.
\par\noindent\textbf{Вывод:} Сумма приведённых теплот за полный обратимый цикл Карно равна нулю. Этот результат обобщается на любой обратимый цикл (теорема Клаузиуса): $\oint \frac{\delta Q_{\text{обр}}}{T} = 0$.
\end{theorem}

\begin{definition}[Термодинамическое определение энтропии]
Поскольку интеграл от приведённой теплоты по любому замкнутому обратимому пути равен нулю, величина $\frac{\delta Q_{\text{обр}}}{T}$ является полным дифференциалом некоторой функции состояния системы $S$, называемой \textbf{энтропией}:
\begin{equation}
    dS = \frac{\delta Q_{\text{обр}}}{T}. \label{eq:entropy_def}
\end{equation}
Изменение энтропии при переходе системы из состояния 1 в состояние 2 не зависит от пути процесса и вычисляется по формуле:
\begin{equation}
    \Delta S = S_2 - S_1 = \int_1^2 \frac{\delta Q_{\text{обр}}}{T}. \label{eq:delta_S}
\end{equation}
Единица измерения энтропии в СИ --- джоуль на кельвин (Дж/К).
\end{definition}

\begin{theorem}[Математическая запись второго начала термодинамики]
Для произвольного (обратимого или необратимого) процесса изменение энтропии удовлетворяет неравенству Клаузиуса:
\begin{equation}
    dS \ge \frac{\delta Q}{T}, \label{eq:second_law_math}
\end{equation}
где знак ``$=$'' относится к обратимым процессам, а ``$>$'' --- к необратимым.
\par\noindent Для \textbf{изолированной системы} ($\delta Q = 0$) второе начало принимает вид:
\begin{equation}
    dS \ge 0 \quad \text{или} \quad \Delta S \ge 0. \label{eq:second_law_isolated}
\end{equation}
Это означает, что в изолированной системе энтропия либо остаётся постоянной (обратимые равновесные процессы), либо возрастает (необратимые реальные процессы). Убывание энтропии в замкнутой системе невозможно.
\end{theorem}

\begin{remark}
\textbf{Физический смысл и следствия:}
\begin{enumerate}
    \item Энтропия является мерой необратимости рассеяния энергии. Рост энтропии указывает на направление естественных процессов (от порядка к хаосу, от неравновесия к равновесию).
    \item Второе начало термодинамики запрещает процессы, в которых тепло самопроизвольно переходит от холодного тела к горячему без компенсации, или КПД тепловой машины превышает $\eta_{\text{К}}$.
    \item Для неизолированных систем энтропия может уменьшаться ($\Delta S < 0$), но только за счёт передачи тепла окружающей среде, при этом суммарная энтропия ``система + окружающая среда'' всегда возрастает.
\end{enumerate}
\end{remark}
